\begin{abstract}
Between keeping every shot as a noisy training example and averaging all shots
into one clean feature lies an entire spectrum of data organizations that
quantum machine learning usually leaves implicit. We make that choice explicit.
Shot grouping partitions a fixed measurement record into disjoint groups and
averages within each group, tuning smoothly between the two conventions with a
single validated parameter and no additional quantum executions. We evaluate
the method on chaotic synthetic benchmarks and on real-world data under strict
execution budgets. An intermediate grouping consistently outperforms the
standard expectation-value pipeline. Matched controls show that the gain acts
as correctly scaled regularization of the classical readout. Unmitigated runs
on two superconducting processors confirm the trend on every task. Shot
grouping therefore offers a practical way to improve finite-shot quantum
learning at no added quantum cost.
\end{abstract}

\begin{IEEEkeywords}
quantum reservoir computing, finite shot learning, measurement aggregation,
time series forecasting
\end{IEEEkeywords}
